\section{Rise blocks, C3 blocks, and the block-layer balance}

This chapter records the algebra of rise steps, C3 steps, and the
block-layer balance that produces the valuation lower bounds.

\subsection{Rise steps}

A rise step has weight $t=1$ or $t=2$.

For $t=1$,
\begin{equation*}
  2y=5x+1,
\end{equation*}
so modulo $5$,
\begin{equation*}
  2y\equiv1\pmod5,
  \qquad
  y\equiv3\pmod5.
\end{equation*}
For $t=2$,
\begin{equation*}
  4y=5x+1,
\end{equation*}
so
\begin{equation*}
  4y\equiv1\pmod5,
  \qquad
  y\equiv4\pmod5.
\end{equation*}
These are the two residue rules recorded by
\texttt{rise\_\allowbreak step\_\allowbreak next\_\allowbreak
mod\_\allowbreak five}.

\subsection{C3 steps}

A C3 step has weight $t\ge3$.  For every positive state $x$,
\begin{equation*}
  \frac{5x+1}{2^t}\le\frac{5x+1}{8}<x,
\end{equation*}
because $5x+1<8x$ for $x\ge1$.  Hence every C3 step strictly
decreases the state.  This is the elementary fact
\texttt{c3Step\_\allowbreak lt\_\allowbreak of\_\allowbreak pos}.

\subsection{Small block molecules}

For rise blocks, the molecule formulas are:
\begin{center}
\small
\begin{tabular}{@{}lll@{}}
\toprule
block word & molecule $B$ & total weight $U$ \\
\midrule
$[1]$ & $1$ & $1$ \\
$[2]$ & $1$ & $2$ \\
$[1,1]$ & $7$ & $2$ \\
$[1,2]$ & $7$ & $3$ \\
$[2,1]$ & $9$ & $3$ \\
$[2,2]$ & $9$ & $4$ \\
\bottomrule
\end{tabular}
\end{center}
For example,
\begin{equation*}
  \wordA([1,2])=5+2=7,
  \qquad
  \wordA([2,1])=5+4=9.
\end{equation*}

\subsection{A C3 block}

The word $[3,1]$ has molecule
\begin{equation*}
  \wordA([3,1])=5+8=13,
  \qquad
  U=4.
\end{equation*}
Its margins were analysed in the PMI-B chapter: the proper prefix of
length $1$ is bad.

\subsection{Block equation}

For any block with weights $u_1,\dots,u_L$ and total weight $U$,
\begin{equation*}
  2^U\,r_L=5^L\,r_0+B.
\end{equation*}
This is the exact word-orbit identity restricted to the block.

\subsection{The reset block}

The reset block starts at the candidate predecessor $x$ and maps it
to the block head $r_j$ with weight $t\in\{1,2\}$:
\begin{equation*}
  2^t\,r_j=5x+1.
\end{equation*}
With
\begin{equation*}
  x=5^{k_0}s+\delta5^{j-1}-1,
\end{equation*}
this becomes
\begin{equation*}
  t=2:\quad 4(r_j+1)=5^{k_0+1}s+\delta5^j,
\end{equation*}
\begin{equation*}
  t=1:\quad 2(r_j+1)=5^{k_0+1}s+5^j-2.
\end{equation*}
These are the reset equations used by the corrected window bound.

\subsection{The block-layer balance}

The compiled lemma \texttt{rise\_\allowbreak block\_\allowbreak balance}
derives the valuation lower bounds from the block structure:
\begin{equation*}
  t=2:\quad
  \vtwo{5^{k_0+1}s+\delta5^j}\ge2j+11,
\end{equation*}
\begin{equation*}
  t=1:\quad
  \vtwo{5^{k_0+1}s+5^j-2}\ge2j+12.
\end{equation*}
The lower bounds come from counting the rise and C3 steps in the
block and from the exact divisibility of the reset equation.  The
paper treats the lemma as compiled input; it does not re-derive it in
the main text.

\subsection{Why the lower bounds are sharp}

The arithmetic counterexample
\begin{equation*}
  (j,k_0,t,\delta,s)=(36,0,2,3,2^{83}-3\cdot5^{35})
\end{equation*}
attains valuation $83=2\cdot36+11$, exactly the lower bound of the
$t=2$ balance.  The corrected window bound would require at most
$82$.  This is the sharpness point: the block-layer lower bound and
the reachable upper bound are separated by exactly one.

\subsection{Status}

\begin{center}
\small
\begin{tabular}{@{}p{0.56\textwidth}p{0.36\textwidth}@{}}
\toprule
Statement & Status \\
\midrule
rise residue rules & compiled \\
C3 decrease fact & compiled \\
block molecule formulas & math proven \\
reset equations & compiled \\
\texttt{rise\_\allowbreak block\_\allowbreak balance} & compiled \\
corrected window bound & compiled (conditional) \\
\bottomrule
\end{tabular}
\end{center}
